\documentclass{tietreport}

\usepackage{graphicx}
\usepackage{float}
\usepackage{booktabs}
\usepackage{tabularx}
\usepackage{array}
\usepackage{enumitem}
\usepackage{listings}
\usepackage{xcolor}
\usepackage[
  date=year,
  style=alphabetic,
  backend=biber,
  sorting=nty,
  maxnames=5,
  minnames=3,
  maxalphanames=4,
  minalphanames=3,
  backref=true,
  doi=false,
  isbn=false,
  url=true,
  eprint=false
]{biblatex}
\addbibresource{references.bib}

\lstset{
  basicstyle=\ttfamily\small,
  breaklines=true,
  frame=single,
  columns=fullflexible,
  keepspaces=true,
  showstringspaces=false
}
\setlist[itemize]{leftmargin=1.5em,itemsep=3pt,topsep=3pt}
\setlist[enumerate]{leftmargin=1.7em,itemsep=3pt,topsep=3pt}
\renewcommand{\arraystretch}{1.2}

\TitlePageHeader{UCS503P -- Software Engineering Project}
\ProjectTitle{LedgerX -- Web3 Wallet Protection System}
\ProjectSubTitle{A Decentralized Asset Protection Platform with Transaction Anomaly Detection}
\author{%
  \texttt{1024030433} Ishan Rehan\\
  \texttt{1024030812} A S Danish\\
  \texttt{1024030813} Shreyansh Dua\\
  \texttt{1024030483} Shaurya Garg%
}
\TitlePageSubText{BE Computer Science and Engineering}
\AdvisorName{Ms. Nisha Thakur}
\TitlePageFooterText{{\large Project Report}\\[0.2cm]
  \bfseries Computer Science and Engineering Department\\[0.15cm]
  Thapar Institute of Engineering and Technology, Patiala%
}

\begin{document}
\maketitle
\tableofcontents
\listoffigures
\listoftables

\chapter{Introduction}

\section{Background and Context}
Blockchain platforms allow users to control digital assets through cryptographic
wallets and smart contracts. Transactions are recorded on a distributed ledger
and are generally difficult to reverse after confirmation. This property supports
transparent settlement, but it also means that a mistaken or malicious transfer
can cause permanent loss. Users therefore need tools that help them inspect
wallet activity, recognize unusual transaction behaviour, and take protective
action before or after a suspicious event.

LedgerX is a Web3 wallet and asset-management project that combines a browser
application, Ethereum-compatible smart contracts, and a Python-based anomaly
detection service. The repository includes a Next.js client, Solidity contracts
developed with Hardhat, and a FastAPI machine-learning service. The platform
also contains token creation, token listing, liquidity-pool interaction,
transaction history, and transfer-related functionality.

\section{Problem Statement}
A conventional wallet primarily provides signing and transaction submission.
It does not necessarily explain whether a transaction is unusual relative to
the wallet's historical behaviour. A separate monitoring and analysis layer
can provide additional context. LedgerX addresses this problem by organizing
wallet operations, transaction monitoring, anomaly scoring, and emergency
asset-exit mechanisms into one project.

The project should not be interpreted as a guarantee that all attacks can be
prevented or that every asset can always be recovered. Detection depends on
available transaction data and model behaviour; on-chain actions depend on
contract permissions, liquidity, network conditions, and successful execution.

\section{Project Objectives}
The objectives of LedgerX are:
\begin{enumerate}
  \item Provide a web interface for wallet connection and common wallet/token
  operations.
  \item Implement Solidity contracts for token creation, wallet custody,
  token transfer, liquidity management, and guarded emergency exit.
  \item Build a transaction anomaly detection pipeline using Python data
  processing and an Isolation Forest model.
  \item Expose anomaly detection functionality through a FastAPI service.
  \item Provide transaction and token views so users can inspect activity.
  \item Document the architecture, test strategy, limitations, and deployment
  requirements of the integrated project.
\end{enumerate}

\section{Scope}
The delivered repository contains the frontend, contract, and model components.
The report distinguishes implemented modules from the intended end-to-end
workflow. In particular, the diagrams supplied with the project illustrate the
proposed full protection journey; they should not be read as proof that every
step is automatically connected in production. The machine-learning service
returns anomaly analysis, while automatic invocation of the Solidity emergency
exit requires a secure integration, authorization policy, and operational
monitoring layer.

\chapter{Methodology and Design}

\section{Development Methodology}
The project is divided into three main engineering areas: client application,
smart-contract layer, and anomaly-detection service. This modular structure
allows each area to be developed and tested independently before integration.
The frontend communicates with wallet providers and application APIs. The
contract layer implements asset custody and exchange-related operations.
The model service processes transaction features and returns anomaly results.

\section{System Architecture}
\subsection{High-Level Design}
The client is implemented using Next.js and React. Wallet connectivity is
configured using Wagmi and RainbowKit, while Ethers is included for contract
interaction. Server-side API routes handle operations such as anomaly requests,
crypto-related requests, AI-assisted transfer parsing, and token metadata
persistence. MongoDB is used by the token metadata API.

The model service is implemented in Python. Its pipeline fetches or accepts
transaction data, cleans and transforms the data, applies an Isolation Forest
detector, and serializes the results. FastAPI exposes service endpoints.
The smart-contract project uses Solidity and Hardhat, with OpenZeppelin
utilities for safer token transfers and reentrancy protection.

\begin{figure}[H]
\centering
\includegraphics[width=\textwidth,height=0.72\textheight,keepaspectratio]{deployment.png}
\caption{LedgerX deployment architecture supplied with the project}
\label{fig:deployment}
\end{figure}

\subsection{Component Architecture}
Figure~\ref{fig:components} shows the major software components and their
relationships. The user-facing pages call reusable wallet and API utilities.
The anomaly service is a separate Python component, and the Solidity contracts
are deployed to an EVM-compatible network.

\begin{figure}[H]
\centering
\includegraphics[width=\textwidth,height=0.72\textheight,keepaspectratio]{component.png}
\caption{LedgerX component diagram}
\label{fig:components}
\end{figure}

\subsection{Data Flow}
The data-flow diagram summarizes the movement of user and transaction
information between the wallet, analysis service, and supporting data stores.
Transaction details are transformed into features before they are evaluated
by the anomaly detector. The resulting classification can be presented to
the user or used by a separately authorized protection workflow.

\begin{figure}[H]
\centering
\includegraphics[width=\textwidth,height=0.68\textheight,keepaspectratio]{dfd.png}
\caption{LedgerX data-flow diagram}
\label{fig:dfd}
\end{figure}

\section{Technical Stack}
\begin{table}[H]
\centering
\begin{tabularx}{\textwidth}{>{\bfseries}p{0.23\textwidth}X}
\toprule
Layer & Technologies observed in the repository\\
\midrule
Frontend & Next.js 14, React 18, TypeScript, CSS, Tailwind-related tooling\\
Wallet integration & Wagmi, RainbowKit, Viem, Ethers\\
Backend APIs & Next.js route handlers and FastAPI\\
Smart contracts & Solidity, Hardhat, OpenZeppelin contracts\\
Anomaly detection & Python, Pandas, NumPy, scikit-learn, Isolation Forest\\
Persistence & MongoDB client for token metadata\\
External data & Etherscan API integration and cryptocurrency price API integration\\
Development & npm, Node.js 22 target, Python environment, Git\\
\bottomrule
\end{tabularx}
\caption{Technology stack}
\label{tab:stack}
\end{table}

\section{Use-Case Model}
The primary actor is the wallet owner, who can connect a wallet, create or
manage protected wallet functionality, transfer funds, monitor activity, and
view transactions. An attacker is represented as an external actor in the
design diagram to model unauthorized access attempts. Blockchain services
participate in token and liquidity operations.

\begin{figure}[H]
\centering
\includegraphics[width=\textwidth,height=0.75\textheight,keepaspectratio]{usecase.png}
\caption{LedgerX use-case diagram}
\label{fig:usecase}
\end{figure}

\section{Workflow Models}
The activity and sequence diagrams show the intended protection workflow:
wallet connection, protected wallet setup, funding, monitoring, threat
assessment, emergency exit, and transaction review. The no-threat branch
returns to monitoring. The threat branch proceeds to the emergency action
and liquidity-pool interaction as represented in the supplied design.

\begin{figure}[H]
\centering
\includegraphics[width=0.70\textwidth,height=0.78\textheight,keepaspectratio]{activity.png}
\caption{LedgerX activity diagram}
\label{fig:activity}
\end{figure}

\begin{figure}[H]
\centering
\includegraphics[width=\textwidth,height=0.68\textheight,keepaspectratio]{sequential.png}
\caption{LedgerX sequence diagram}
\label{fig:sequence}
\end{figure}

\begin{figure}[H]
\centering
\includegraphics[width=\textwidth,height=0.68\textheight,keepaspectratio]{swimlane.png}
\caption{LedgerX swimlane workflow}
\label{fig:swimlane}
\end{figure}

\section{Class and Contract Model}
The class diagram supplied for the proposal identifies the conceptual
relationship between an owner, wallet, anomaly guard, and liquidity pool.
The repository implements these responsibilities through separate Solidity
contracts rather than one monolithic contract.

\begin{figure}[H]
\centering
\includegraphics[width=0.88\textwidth,height=0.68\textheight,keepaspectratio]{classdiagram.png}
\caption{Conceptual class diagram}
\label{fig:class}
\end{figure}

\section{Implementation Details}

\subsection{Frontend and Wallet Interface}
The client is organized as a Next.js application with pages for the dashboard,
wallet, token creation, token listing, transfers, transactions, liquidity, and
graph-oriented views. Reusable components include navigation, wallet
connection, transaction listing, and trading-related interface components.
Wallet configuration is defined in the Wagmi setup, which includes a local
Hardhat chain and supported public networks. The application reads contract
addresses from environment variables and supplies local development defaults.

\subsection{Wallet and Token Smart Contracts}
The Solidity project includes the following contracts:
\begin{itemize}
  \item \textbf{Wallet:} maintains per-user ERC-20 token balances. Deposits
  transfer tokens into contract custody using \texttt{SafeERC20}; withdrawals
  transfer tokens back to the caller; internal transfers update the ledger.
  \item \textbf{CustomToken:} defines the ERC-20 token implementation used by
  the token creation flow.
  \item \textbf{TokenFactory:} deploys a new custom token and emits a
  \texttt{TokenCreated} event containing token metadata and creator details.
  \item \textbf{TransferTo:} transfers ERC-20 tokens from the caller to a
  recipient using \texttt{safeTransferFrom}.
  \item \textbf{LiquidityPool:} implements token/ETH liquidity provision and
  swaps using a constant-product pricing formula with a 0.3\% fee.
  \item \textbf{AnomalyGuardWallet:} provides an owner-controlled emergency
  exit. It attempts to swap held tokens through the configured pool and, if
  the swap fails, transfers the tokens to the owner as a fallback.
\end{itemize}

\subsection{Liquidity Pool Design}
The liquidity pool maintains explicit token and ETH reserves. Liquidity
providers receive shares proportional to the pool contribution. A minimum
liquidity amount is locked on initial provisioning. Swaps use the
constant-product formula:
\[
  \mathrm{amountOut} =
  \frac{\mathrm{amountInWithFee}\cdot \mathrm{reserveOut}}
       {\mathrm{reserveIn}\cdot 1000+\mathrm{amountInWithFee}},
  \qquad
  \mathrm{amountInWithFee}=997\cdot\mathrm{amountIn}.
\]
The contract uses slippage bounds and a reentrancy guard. These mechanisms
reduce specific classes of implementation risk but do not eliminate all
economic or smart-contract risks.

\subsection{Anomaly Guard and Emergency Exit}
The guard contract stores an owner, an ERC-20 token address, and a liquidity
pool address. The owner can forward tokens or ETH into the contract. The
emergency exit attempts to convert the token balance to ETH through the pool.
If the swap reverts, the fallback path returns the token balance to the owner.
This is an owner-triggered contract function in the inspected implementation;
it is not, by itself, an autonomous on-chain detector.

\subsection{Machine-Learning Pipeline}
The model service is structured as a sequence of data acquisition, cleaning,
transformation, model training, model persistence, and anomaly inference.
The main orchestration module uses an Etherscan API helper, data-cleaning and
data-transformation classes, and an Isolation Forest detector. The repository
also contains serialized model artifacts and example transaction/result data.

Isolation Forest is an unsupervised anomaly-detection method that isolates
observations using randomized partitioning. Observations that are isolated
with shorter paths are more likely to be anomalous. The method is suitable
for exploratory transaction screening because it does not require a complete
set of labelled attack examples \cite{liu2008isolation,liu2012isolation}.
Its output is dependent on feature selection, training data, contamination
settings, and threshold choices.

\subsection{FastAPI Service}
The model API defines transaction input structures with fields such as hash,
timestamp, value, gas, and gas price. It provides detection and model-training
functionality. Training is guarded by a server-side key; if the key is not
configured, training is disabled. The API limits transaction-list size and
configures CORS through environment settings. These controls help constrain
requests to the service.

\subsection{Application API Routes and Persistence}
The Next.js application contains API routes for anomaly requests, crypto
information, AI-assisted transfer parsing, and token metadata storage. The
token-saving route validates token addresses and bounds metadata fields before
using MongoDB. Optional image uploads are routed to ImgBB when an API key is
configured. External services therefore require correctly configured secrets
and network access.

\section{Security Considerations}
The project uses OpenZeppelin's \texttt{SafeERC20} and
\texttt{ReentrancyGuard} in relevant contracts. Input validation is present in
selected API routes. Environment variables are used for API credentials,
database connection strings, and deployment addresses. The deployment
documentation also identifies environment-specific build requirements.

Security review should additionally consider owner-key compromise, malicious
tokens, price manipulation, liquidity exhaustion, front-running, oracle
integrity, model evasion, training-data poisoning, API abuse, and incorrect
contract address configuration. A production deployment should use
independent smart-contract auditing, restricted administrative access,
monitoring, and a carefully tested incident-response policy.

\chapter{Results and Evaluation}

\section{Implementation Inventory}
The repository contains the following major implementation artifacts:
\begin{table}[H]
\centering
\begin{tabularx}{\textwidth}{p{0.28\textwidth}X}
\toprule
Area & Evidence in the supplied repository\\
\midrule
Frontend & Next.js pages, wallet components, hooks, API route handlers\\
Contracts & Wallet, CustomToken, TokenFactory, TransferTo, LiquidityPool,
AnomalyGuardWallet\\
Model service & FastAPI app, model orchestration, Isolation Forest detector,
data processing modules\\
Model artifacts & Serialized Isolation Forest, scaler, thresholds, example
transaction and result JSON files\\
Tests & Python tests for detection, API, data cleaning, integration, and
visualization; a Hardhat liquidity-pool test file\\
Deployment & Deployment notes and Hardhat Ignition deployment artifacts\\
\bottomrule
\end{tabularx}
\caption{Implementation inventory based on the project ZIP}
\label{tab:inventory}
\end{table}

\section{Testing Methodology}
The repository includes Python test modules covering anomaly detection,
API behaviour, ARIMA functionality, data cleaning, detection quality,
integration, and visualization. The Solidity project includes a test file
for the liquidity pool. These files provide a test foundation; the presence
of test files alone does not establish that every test passes in the current
environment.

A complete verification run should include:
\begin{enumerate}
  \item Python unit tests for cleaning, transformation, model inference, and
  API validation.
  \item Hardhat tests for deposits, withdrawals, swaps, slippage, reserve
  accounting, and emergency-exit fallback behaviour.
  \item Frontend production build and manual checks of wallet connection,
  network selection, API error states, and transaction views.
  \item Integration tests using a local Hardhat network and controlled
  transaction examples.
\end{enumerate}

\section{Example Anomaly Output}
The supplied \texttt{results.json} contains example records illustrating a
normal transaction and an anomalous transaction. The anomalous example is
flagged with high-value transaction, high gas consumption, high gas price,
and unusual value/gas ratio indicators. These are example output records and
are not a benchmark report or a measured accuracy evaluation.

\begin{table}[H]
\centering
\begin{tabularx}{\textwidth}{p{0.30\textwidth}XX}
\toprule
Example record & Classification & Reported indicators\\
\midrule
Transaction \texttt{0x123...} & Normal & No anomalies detected\\
Transaction \texttt{0x456...} & Anomalous & High value, high gas usage,
high gas price, unusual value/gas ratio\\
\bottomrule
\end{tabularx}
\caption{Illustrative records in the supplied model results}
\label{tab:examples}
\end{table}

\section{Performance Metrics}
No reproducible latency, throughput, precision, recall, F1-score, or
false-positive benchmark is asserted in this report because the supplied
materials do not provide a documented evaluation run with a defined test
dataset and measurement procedure. These metrics should be measured before
making quantitative performance claims.

Recommended evaluation measures include:
\begin{itemize}
  \item \textbf{Precision:} fraction of flagged transactions that are truly
  anomalous in a labelled evaluation set.
  \item \textbf{Recall:} fraction of known anomalous transactions detected.
  \item \textbf{False-positive rate:} fraction of normal transactions
  incorrectly flagged.
  \item \textbf{Inference latency:} elapsed time from API request receipt to
  anomaly response.
  \item \textbf{Contract execution success:} proportion of controlled
  emergency-exit scenarios completing the intended on-chain action.
\end{itemize}

\section{Analysis and Current Limitations}
The project demonstrates a multi-layer implementation spanning frontend,
backend APIs, machine-learning analysis, and Solidity contracts. Its modular
design supports independent development and future integration. The key
limitation is that anomaly classification and on-chain emergency exit are
separate capabilities unless an authenticated service or user-authorized
workflow connects them. A machine-learning result cannot directly execute
a smart contract without a transaction signer, keeper/oracle mechanism, or
explicit user action.

Other limitations include dependence on external APIs and environment
configuration, the need for robust labelled evaluation data, the risk of
false alarms, and the economic constraints of liquidity-based recovery.
The fallback path in the guard contract returns tokens if the pool swap
fails; it does not guarantee conversion into a stable asset.

\chapter{Conclusion}

\section{Summary of Work}
LedgerX is a Web3 asset-management and monitoring project combining a
Next.js/React interface, Solidity smart contracts, and a Python anomaly
detection service. The repository includes wallet and token operations,
liquidity-pool functionality, transaction-related interfaces, a FastAPI
service, and an Isolation Forest-based detection pipeline. The supplied
design diagrams document the intended wallet-protection workflow and
system structure.

\section{Key Findings}
\begin{itemize}
  \item A modular architecture separates user interaction, blockchain
  execution, and transaction analysis.
  \item The Solidity implementation includes guarded token custody,
  constant-product pool swaps, slippage checks, and an owner-controlled
  emergency exit.
  \item The anomaly-detection pipeline provides a foundation for screening
  transaction behaviour, but its effectiveness must be established through
  labelled and reproducible evaluation.
  \item Automatic threat detection and contract execution require a secure
  integration layer and should not be assumed solely from the presence of
  the individual modules.
\end{itemize}

\section{Future Work and Recommendations}
\begin{enumerate}
  \item Integrate model alerts with a secure, auditable transaction
  authorization workflow.
  \item Add labelled datasets and report precision, recall, F1-score,
  false-positive rate, and inference latency.
  \item Extend contract tests to cover adversarial tokens, reentrancy,
  slippage, pool depletion, and emergency-exit edge cases.
  \item Add independent smart-contract audits and formal verification for
  critical asset-handling functions.
  \item Improve wallet-level monitoring, alert delivery, and transaction
  explanations.
  \item Add deployment automation, secrets management, structured logs,
  and health monitoring for all services.
\end{enumerate}

\section{Final Remarks}
LedgerX demonstrates how Web3 interfaces, smart contracts, and machine
learning can be brought together in a software-engineering project focused
on digital-asset safety. The work provides a foundation for further
integration and evaluation. Responsible deployment requires careful
security review, transparent limitations, and measured evidence of
detection and recovery performance.

\printbibliography[heading=bibintoc]

\end{document}
